\clearpage
\fancyhead[L]{\small BRAMBILLA. EFFEMERIDI ASTRONOMICHE.}
\SourceLeaf{634}
\begin{center}\small Göttingische gelehrte Anzeigen. Stück 64. Seiten 633..637. 21. April 1823.\\[3pt]\rule{.72\linewidth}{.3pt}\end{center}

\emph{Effemeridi Astronomiche di Milano per l'anno \textup{1823} calcolate da} \textsc{Enrico Brambilla}. \emph{Con Appendice.} Mailand 1822. Die Ephemeriden selbst umfassen 112 Seiten, der Anhang 88 Seiten im Oktavformat.

Man ist von den Mailänder Ephemeriden seit Langem gewohnt, dass sie einige Aufsätze enthalten. Deren Wert reicht über das Erscheinungsjahr hinaus, und sie machen die Ephemeriden auch für Menschen interessant, die keinen Anlass haben, den astronomischen Kalender zu benutzen.
Über den Kalender brauchen wir nur zu sagen, dass seine gelungene Anlage unverändert geblieben ist. Der Anhang enthält die folgenden Aufsätze.

Geografische Lage einiger von Mailand aus sichtbarer Berge von \textsc{Barnabas Oriani}.
Dieser hochverdiente Astronom teilt uns hier ein Verzeichnis von 49 Punkten in Oberitalien und der Schweiz mit, das ihre Breite und Länge, ihre Lage in Bezug auf die Kathedrale von Mailand und ihre Höhe über dem Meeresspiegel enthält.
Diese Bestimmungen beruhen größtenteils auf den Messungen, die die Mailänder Astronomen in den Jahren 1788---1791 auf Anordnung der österreichischen und in den Jahren 1803---1806 auf Anordnung der italienischen Regierung ausgeführt haben.
Diese Messungen erstreckten sich nicht bis zum Meer, sodass nur relative Höhen bestimmt wurden. Ihrer Reduktion auf den Meeresspiegel wurde die barometrische Höhenbestimmung von Mailand zugrunde gelegt.
Einige Bestimmungen des Verzeichnisses beruhen auf Messungen aus jüngster Zeit, die hier etwas ausführlicher beschrieben werden und von besonderem Interesse sind.
Die Lage des Monte Viso im Piemont wurde durch Messungen von Turin und Mailand aus bestimmt. Von Mailand ist er 96380 Toisen (25 geografische Meilen) entfernt.
Die von beiden Orten aus vorgenommenen Höhenmessungen stimmen am besten überein, wenn man für die terrestrische Refraktion das 0.08-Fache der Krümmung des terrestrischen Bogens ansetzt. \textsc{Oriani} setzt daher die Höhe über dem Meeresspiegel mit 1968 Toisen an.
Die Bestimmung der Höhe des Monte Cimone im Apennin unweit von Lucca auf 1112 Toisen ist besonders bemerkenswert, da sie auf der unmittelbaren Messung der Depression des Meereshorizonts sowohl des Mittelmeers als auch der Adria beruht. \textsc{Brioschi} hat diese Messungen im Jahr 1817 ausgeführt.
Als Höhen werden für den Monte Rosa 2385, für das Finsterarhorn 2203 und für den Simplon 1805 Toisen angesetzt. --- Bericht über die 1822 ausgeführten Operationen zur Bestimmung der Längenunterschiede zwischen mehreren Orten in Italien durch Pulversignale auf dem Monte Cimone, von \textsc{Francesco Carlini}.
Durch die Triangulierungen in Frankreich und den österreichischen Staaten wird man, sobald sie abgeschlossen und miteinander verknüpft sind, die Messung eines überaus großen Bogens des mittleren Parallelkreises erhalten, der vom Atlantischen Meer bis Orsova etwa 24 Grad betragen wird.
Es liegt auf der Hand, wie wichtig diese Messungen in Verbindung mit geeigneten astronomischen Operationen für eine vollständigere Kenntnis der Gestalt der Erde werden können.
Die Beschaffenheit der Landstriche selbst, durch die dieser Bogen verläuft, ist für die Operationen zur Bestimmung des Längenunterschieds der Endpunkte besonders günstig. Auf dieser Strecke liegen so viele hohe Berge mit einer ungeheuer weiten Aussicht, dass man mit einer verhältnismäßig sehr kleinen Zahl von Zwischenpunkten auskommen kann.
Ein erster Versuch dieser Art wurde bereits im Sept. 1821 unternommen, als man den Längenunterschied zwischen der Mailänder Sternwarte und dem Hospiz auf dem Mont Cenis durch Pulversignale auf der 1792 Toisen hohen und 86000 Toisen von Mailand entfernten Rocca Melone bestimmte.
Durch den glücklichen Ausgang dieses Versuchs ermutigt, wollte man zu einer
\SourceLeaf{635} umfassenderen Verbindung übergehen.
Dazu schlugen die französischen Geografen einen kühnen Plan vor. Man meinte, mithilfe dreier Zwischenpunkte, nämlich des oben erwähnten Monte Viso, des Monte Cero bei Padua und des Monte Maggiore im Friaul, in Einer Nacht die Verbindung zwischen der Ostküste des Adriatischen Meeres und dem Mont d'Or bei Clermont mitten in Frankreich herstellen zu können.
Man nahm jedoch wieder davon Abstand, weil man die Schwierigkeiten für zu groß hielt.
Es ist noch ungewiss, ob der Monte Viso überhaupt bestiegen werden kann, und noch ungewisser, ob man sich während der Nacht auf seinem höchsten, steilen Gipfel aufhalten kann.
Und selbst wenn sich diese Schwierigkeiten überwinden ließen, fürchtet man, dass das Licht der Pulverblitze bei der ungeheuren Entfernung vom Monte Cero (50 geogr. Meilen) selbst durch die stärksten Fernrohre nicht zu sehen wäre
(Nach diesen Äußerungen scheint dieser riesenhafte Plan noch nicht endgültig aufgegeben zu sein. Der Rezensent stellt jedoch anhand der Angaben zu den Höhen der beiden Punkte und zu ihrer Entfernung fest, dass keiner der beiden über dem physischen Horizont des anderen liegen kann. Ohne diesen Umstand ließe sich den beiden letzten Schwierigkeiten durch den Einsatz großer Heliotrope begegnen).
Man entschied sich daher vorerst für eine begrenztere Operation und ließ Anfang Mai 1822 mehrere Nächte lang auf dem Monte Cimone Pulversignale geben, die auf dem Monte Cero, in Mailand und an verschiedenen anderen italienischen Sternwarten beobachtet werden sollten.
Doch das ungünstige Wetter verhinderte, dass diese Operationen in der Hauptsache Erfolg hatten. Weder in Mailand noch auf dem Monte Cero wurden die Pulversignale gesehen.
In Parma, Modena, Bologna und Florenz wurden sie hingegen beobachtet. Doch die zum Teil beträchtlichen Abweichungen der Ergebnisse von den Längendifferenzen, die sich aus den geodätischen Operationen ergeben hatten, scheinen zu beweisen, dass die Zeitbestimmung nicht an allen diesen Orten die nötige Genauigkeit hatte.
Ganz besonders bemerkenswert ist noch, dass Herr \textsc{Carlini} aus seinen Beobachtungen mit einem 18-zölligen \textsc{Reichenbach}'schen Repetitionskreis für Parma eine um $22''6$ größere Polhöhe ermittelt hat, als sich aus der geodätischen Verbindung mit Mailand ergab.
Da sich ähnliche Anomalien schon früher an mehreren anderen Orten Oberitaliens gezeigt haben, fällt es schwer, an ihrer Realität zu zweifeln, und es ist sehr zu wünschen, dass alle Hauptdreiecke der österreichischen Triangulierungen bald vollständig bekannt gemacht werden. --- Die übrigen Beiträge des Anhangs enthalten außerdem die von \textsc{Angelo Cesaris} in den Jahren 1817 und 1818 beobachteten Oppositionen des Uranus sowie die von demselben Astronomen im Jahr 1821 beobachteten Oppositionen des Jupiter und des Saturn. Hinzu kommen die von \textsc{Hallaschka} in Prag beobachteten Sternbedeckungen und Verfinsterungen der Jupitermonde sowie seine Beobachtungen des ersten Kometen von 1822. Den Schluss bilden die meteorologischen Beobachtungen in Mailand aus dem Jahr 1820 von A. \textsc{Cesaris}.
